# ============================================================
# Model Fitting: ARIMA and Dynamic Regression (Phillips Curve)
# ============================================================
#
# Structure:
#   1.  Data loading and preparation
#   2.  Structural break dummies
#   3.  Train / test split  (last 24 months = test)
#   4.  ARIMA baseline
#       4.1  Auto ARIMA
#       4.2  Manual candidate grid
#       4.3  ARIMA with structural break dummies
#       4.4  Residual diagnostics and Ljung-Box
#       4.5  Forecast and accuracy
#   5.  General-to-Specific (GETS) variable selection
#       5.1  Build General Unrestricted Model (GUM)
#       5.2  GETS reduction with indicator saturation (IIS + SIS)
#       5.3  Inspect retained regressors
#   6.  Dynamic regression models (fpp3)
#       6.1  Model A: DK unemployment only
#       6.2  Model B: DK unemployment + EU27 inflation
#       6.3  Model C: GETS-selected specification
#       6.4  Residual diagnostics
#       6.5  Forecast and accuracy
#   7.  Head-to-head model comparison
#   8.  Time Series Cross-Validation (TSCV) on best model
# ============================================================

rm(list = ls())

library(tidyverse)
library(lubridate)
library(knitr)
library(ggplot2)
library(patchwork)
library(fpp3)
library(tseries)
library(feasts)
library(strucchange)
library(gets)
library(lmtest)
conflicts_prefer(dplyr::lag)


# ============================================================
# 1. Data loading and preparation
# ============================================================

script_dir <- tryCatch(
  dirname(normalizePath(sys.frame(1)$ofile, mustWork = FALSE)),
  error = function(e) getwd()
)
if (!nzchar(script_dir) || is.na(script_dir)) script_dir <- getwd()

geo_label_to_code <- c(
  "Denmark"                                   = "DK",
  "European Union - 27 countries (from 2020)" = "EU27_2020"
)

laes_eurostat_csv <- function(sti) {
  read_csv(sti, show_col_types = FALSE) |>
    filter(geo %in% names(geo_label_to_code)) |>
    transmute(
      geo    = unname(geo_label_to_code[geo]),
      dato   = ym(TIME_PERIOD),
      vaerdi = as.numeric(OBS_VALUE)
    ) |>
    filter(!is.na(dato)) |>
    arrange(geo, dato)
}

inflation   <- laes_eurostat_csv(file.path(script_dir, "prc_hicp_minr__custom_21465689_linear.csv"))
arbejdsloes <- laes_eurostat_csv(file.path(script_dir, "une_rt_m__custom_21465700_linear.csv"))

inflation_wide <- inflation |>
  pivot_wider(names_from = geo, values_from = vaerdi) |>
  rename(dk_inflation = DK, eu_inflation = EU27_2020)

arbejdsloes_wide <- arbejdsloes |>
  pivot_wider(names_from = geo, values_from = vaerdi) |>
  rename(dk_unemployment = DK, eu_unemployment = EU27_2020)

df <- full_join(inflation_wide, arbejdsloes_wide, by = "dato") |>
  arrange(dato) |>
  drop_na()

combined <- df |>
  mutate(yearmonth = yearmonth(dato)) |>
  select(yearmonth, dk_inflation, eu_inflation, dk_unemployment, eu_unemployment) |>
  as_tsibble(index = yearmonth) |>
  drop_na()

# First-difference unemployment series (confirmed I(1) in stationarity tests).
# Inflation series are sourced as YoY rates of change and used in levels (I(0)).
combined <- combined |>
  mutate(
    diff_dk_unemployment = difference(dk_unemployment),
    diff_eu_unemployment = difference(eu_unemployment)
  )

cat("Sample period:", as.character(min(combined$yearmonth)),
    "to", as.character(max(combined$yearmonth)), "\n")
cat("Observations:", nrow(combined), "\n")


# ============================================================
# 2. Structural break dummies
# ============================================================

# QLR test (initial_data_inspection.R) identified two structural breaks in
# Danish inflation:
#   Break 1: December 2012  ->  step up from January 2013
#            (tail end of the European sovereign debt crisis)
#   Break 2: August 2021    ->  step up from September 2021
#            (onset of the post-COVID inflation surge)
#
# Step dummies take the value 0 before the break and 1 from the break onward.
# Including them as regressors controls for permanent shifts in the mean of
# Danish inflation, preventing the ARIMA error structure from absorbing what
# is in reality a deterministic level shift.

combined <- combined |>
  mutate(
    break_2013 = as.numeric(yearmonth >= yearmonth("2013 Jan")),
    break_2021 = as.numeric(yearmonth >= yearmonth("2021 Sep"))
  )

cat("Break dummy sums:", sum(combined$break_2013), "and", sum(combined$break_2021), "\n")


# ============================================================
# 3. Train / test split
# ============================================================

# Test set: last 24 months.
# Training set: everything prior to the test period.

test_end   <- max(combined$yearmonth)
test_start <- test_end - 23

train <- combined |> filter(yearmonth <  test_start)
test  <- combined |> filter(yearmonth >= test_start)

cat("Training period:", as.character(min(train$yearmonth)),
    "to", as.character(max(train$yearmonth)),
    "(", nrow(train), "obs)\n")
cat("Test period    :", as.character(min(test$yearmonth)),
    "to", as.character(max(test$yearmonth)),
    "(", nrow(test), "obs)\n")


# ============================================================
# 4. ARIMA baseline
# ============================================================

# The ARIMA model forecasts Danish inflation using only its own past values.
# It serves as the benchmark: any model with economic content should beat it.


# --- 4.1 Auto ARIMA ---

fit_arima_auto <- train |>
  model(
    arima_auto = ARIMA(dk_inflation, stepwise = FALSE, approximation = FALSE)
  )

report(fit_arima_auto)


# --- 4.2 Manual candidate grid ---

# Candidates informed by the ACF / PACF of dk_inflation from inspection.
# dk_inflation is I(0), so d = 0 is the natural starting point, though
# auto ARIMA may select d = 1 if the unit root test is borderline.
# We estimate both to let AICc decide.

fit_arima_grid <- train |>
  model(
    arima_000  = ARIMA(dk_inflation ~ pdq(0,0,0)),
    arima_100  = ARIMA(dk_inflation ~ pdq(1,0,0)),
    arima_200  = ARIMA(dk_inflation ~ pdq(2,0,0)),
    arima_010  = ARIMA(dk_inflation ~ pdq(0,1,0)),
    arima_110  = ARIMA(dk_inflation ~ pdq(1,1,0)),
    arima_001  = ARIMA(dk_inflation ~ pdq(0,0,1)),
    arima_002  = ARIMA(dk_inflation ~ pdq(0,0,2)),
    arima_101  = ARIMA(dk_inflation ~ pdq(1,0,1)),
    arima_201  = ARIMA(dk_inflation ~ pdq(2,0,1)),
    arima_011  = ARIMA(dk_inflation ~ pdq(0,1,1)),
    arima_111  = ARIMA(dk_inflation ~ pdq(1,1,1)),
    arima_001s = ARIMA(dk_inflation ~ pdq(0,0,1) + PDQ(0,0,1)),
    arima_101s = ARIMA(dk_inflation ~ pdq(1,0,1) + PDQ(0,0,1)),
    arima_011s = ARIMA(dk_inflation ~ pdq(0,1,1) + PDQ(0,0,1)),
    arima_111s = ARIMA(dk_inflation ~ pdq(1,1,1) + PDQ(0,0,1))
  )

fit_arima_grid |>
  glance() |>
  select(.model, AIC, AICc, BIC, sigma2) |>
  arrange(AICc) |>
  kable(digits = 3, align = "c")


# --- 4.3 ARIMA with structural break dummies ---

# Adding the two step dummies as regressors controls for the identified
# mean shifts in Danish inflation. This should improve fit in-sample and
# reduce forecast bias over the test period if the post-2021 level persists.

fit_arima_breaks <- train |>
  model(
    arima_breaks = ARIMA(
      dk_inflation ~ break_2013 + break_2021,
      stepwise      = FALSE,
      approximation = FALSE
    )
  )

report(fit_arima_breaks)


# --- 4.4 Residual diagnostics ---

fit_arima_auto |>
  gg_tsresiduals() +
  labs(title = "ARIMA Baseline: Residual Diagnostics")

fit_arima_breaks |>
  gg_tsresiduals() +
  labs(title = "ARIMA with Break Dummies: Residual Diagnostics")

# Ljung-Box test: dof = number of AR + MA parameters estimated.
ljung_box_table <- function(fit, lag = 24) {
  dof <- fit |> tidy() |>
    filter(str_detect(term, "^ar|^ma|^sar|^sma")) |>
    nrow()
  fit |>
    augment() |>
    features(.innov, ljung_box, dof = dof, lag = lag)
}

ljung_box_table(fit_arima_auto)   |> kable(digits = 3, align = "c")
ljung_box_table(fit_arima_breaks) |> kable(digits = 3, align = "c")


# --- 4.5 Forecast and accuracy ---

fc_arima_auto   <- fit_arima_auto   |> forecast(new_data = test)
fc_arima_breaks <- fit_arima_breaks |> forecast(new_data = test)

fc_arima_auto |>
  autoplot(combined |> filter(yearmonth >= yearmonth("2020 Jan")), level = c(80, 95)) +
  labs(title = "ARIMA Baseline: Forecast vs. Actual Danish Inflation",
       x = "Month", y = "Inflation (YoY, %)") +
  theme(plot.title = element_text(hjust = 0.5))

fc_arima_breaks |>
  autoplot(combined |> filter(yearmonth >= yearmonth("2020 Jan")), level = c(80, 95)) +
  labs(title = "ARIMA with Break Dummies: Forecast vs. Actual Danish Inflation",
       x = "Month", y = "Inflation (YoY, %)") +
  theme(plot.title = element_text(hjust = 0.5))

bind_rows(
  fc_arima_auto   |> accuracy(test),
  fc_arima_breaks |> accuracy(test)
) |>
  select(.model, RMSE, MAE, MAPE, MASE) |>
  arrange(RMSE) |>
  kable(digits = 3, align = "c")


# ============================================================
# 5. General-to-Specific (GETS) variable selection
# ============================================================

# Strategy:
#   Start from a General Unrestricted Model (GUM) that includes:
#     - AR lags of dk_inflation (lags 1:6)
#     - Contemporaneous and lagged diff_dk_unemployment (lags 0:6)
#     - Contemporaneous and lagged eu_inflation (lags 0:6)
#     - Contemporaneous and lagged diff_eu_unemployment (lags 0:6)
#     - The two structural break step dummies
#   Use isat() with:
#     - IIS (Impulse Indicator Saturation): detects individual outliers
#     - SIS (Step Indicator Saturation): detects additional level shifts
#   Retain regressors significant at the 5% level.
#   The retained specification is then carried into the fpp3 dynamic regression.


# --- 5.1 Build GUM regressor matrix ---

# isat() works on plain numeric vectors, so we extract from the training tsibble.
# All regressors must be I(0): unemployment is first-differenced, inflation in levels.

gets_df <- train |>
  as_tibble() |>
  arrange(yearmonth) |>
  filter(!is.na(diff_dk_unemployment), !is.na(diff_eu_unemployment))

y_tr <- gets_df$dk_inflation
n_tr <- length(y_tr)

max_lag <- 6

build_lag_matrix <- function(x, max_lag, prefix) {
  mat <- sapply(0:max_lag, function(k) {
    c(rep(NA_real_, k), x[seq_len(n_tr - k)])
  })
  colnames(mat) <- paste0(prefix, "_l", 0:max_lag)
  mat
}

x_dk_unemp <- build_lag_matrix(gets_df$diff_dk_unemployment, max_lag, "d_dk_unemp")
x_eu_infl  <- build_lag_matrix(gets_df$eu_inflation,         max_lag, "eu_infl")
x_eu_unemp <- build_lag_matrix(gets_df$diff_eu_unemployment, max_lag, "d_eu_unemp")
x_breaks   <- cbind(
  break_2013 = gets_df$break_2013,
  break_2021 = gets_df$break_2021
)

x_gum_full <- cbind(x_dk_unemp, x_eu_infl, x_eu_unemp, x_breaks)

# Drop rows with any NA (introduced by lagging).
keep  <- complete.cases(x_gum_full)
y_gum <- y_tr[keep]
x_gum <- x_gum_full[keep, , drop = FALSE]

cat("GUM: y observations =", length(y_gum),
    "| regressors =", ncol(x_gum), "\n")


# --- 5.2 GETS reduction with IIS and SIS ---

# IIS detects individual outliers (impulse indicators).
# SIS detects permanent level shifts (step indicators).
# Break dummies are already included in mxreg, so SIS here catches any
# additional breaks not identified by the QLR test.
#
# t.pval = 0.05: retain regressors with |t| > 1.96.
# ar = 1:6: include AR lags of dk_inflation in the GUM.

gets_fit <- isat(
  y      = y_gum,
  ar     = 1:6,
  mxreg  = x_gum,
  iis    = TRUE,
  sis    = TRUE,
  t.pval = 0.05
)

print(gets_fit)


# --- 5.3 Inspect retained regressors ---

cat("\nRetained regressors:\n")
print(coef(gets_fit))


# ============================================================
# 6. Dynamic regression models (fpp3)
# ============================================================

# The three models progress from simple to general, using dk_inflation in
# levels with external regressors that are I(0). ARIMA() selects the error
# order automatically via AICc.
#
# Model A: DK unemployment only (strict domestic Phillips curve).
# Model B: DK unemployment + EU27 inflation (ERM II spillover channel).
# Model C: GETS-selected specification (data-driven, accounts for breaks).
#
# All models condition on actual test-period values of the regressors,
# so forecast errors reflect only the inflation model — not regressor
# forecast uncertainty.


# --- 6.1 Add lagged regressors to the tsibble ---

# Update combined with all lags needed across the three models.
# Adjust the lag selection in Model C based on the GETS output above.

combined <- combined |>
  mutate(
    d_dk_unemp_l0 = diff_dk_unemployment,
    d_dk_unemp_l1 = lag(diff_dk_unemployment, 1),
    d_dk_unemp_l2 = lag(diff_dk_unemployment, 2),
    d_dk_unemp_l3 = lag(diff_dk_unemployment, 3),
    d_dk_unemp_l4 = lag(diff_dk_unemployment, 4),
    d_dk_unemp_l5 = lag(diff_dk_unemployment, 5),
    d_dk_unemp_l6 = lag(diff_dk_unemployment, 6),
    eu_infl_l0    = eu_inflation,
    eu_infl_l1    = lag(eu_inflation, 1),
    eu_infl_l2    = lag(eu_inflation, 2),
    eu_infl_l3    = lag(eu_inflation, 3),
    eu_infl_l4    = lag(eu_inflation, 4),
    eu_infl_l5    = lag(eu_inflation, 5),
    eu_infl_l6    = lag(eu_inflation, 6),
    d_eu_unemp_l0 = diff_eu_unemployment,
    d_eu_unemp_l1 = lag(diff_eu_unemployment, 1),
    d_eu_unemp_l2 = lag(diff_eu_unemployment, 2),
    d_eu_unemp_l3 = lag(diff_eu_unemployment, 3)
  )

train <- combined |> filter(yearmonth <  test_start)
test  <- combined |> filter(yearmonth >= test_start)


# --- 6.2 Model A: DK unemployment only ---

fit_A <- train |>
  model(
    model_A = ARIMA(
      dk_inflation ~ break_2013 + break_2021 +
        d_dk_unemp_l0 + d_dk_unemp_l1 + d_dk_unemp_l2 + d_dk_unemp_l3,
      stepwise      = FALSE,
      approximation = FALSE
    )
  )

report(fit_A)


# --- 6.3 Model B: DK unemployment + EU27 inflation ---

fit_B <- train |>
  model(
    model_B = ARIMA(
      dk_inflation ~ break_2013 + break_2021 +
        d_dk_unemp_l0 + d_dk_unemp_l1 + d_dk_unemp_l2 + d_dk_unemp_l3 +
        eu_infl_l0 + eu_infl_l1 + eu_infl_l2 + eu_infl_l3,
      stepwise      = FALSE,
      approximation = FALSE
    )
  )

report(fit_B)


# --- 6.4 Model C: GETS-selected specification ---

# Replace the right-hand side below with the regressors retained by GETS
# in section 5.3. The template includes all candidates; comment out those
# not retained. Keep break_2013 and break_2021 regardless of GETS outcome
# since they are economically motivated and confirmed by the QLR test.

fit_C <- train |>
  model(
    model_C = ARIMA(
      dk_inflation ~ break_2013 + break_2021 +
        eu_infl_l0 + eu_infl_l1 + eu_infl_l6,
      stepwise      = FALSE,
      approximation = FALSE
    )
  )

report(fit_C)


# --- 6.5 Residual diagnostics ---

fit_A |> gg_tsresiduals() + labs(title = "Model A: Residual Diagnostics")
fit_B |> gg_tsresiduals() + labs(title = "Model B: Residual Diagnostics")
fit_C |> gg_tsresiduals() + labs(title = "Model C: Residual Diagnostics")

ljung_box_table(fit_A) |> kable(digits = 3, align = "c")
ljung_box_table(fit_B) |> kable(digits = 3, align = "c")
ljung_box_table(fit_C) |> kable(digits = 3, align = "c")


# --- 6.6 In-sample information criteria comparison ---

bind_rows(
  glance(fit_arima_auto),
  glance(fit_arima_breaks),
  glance(fit_A),
  glance(fit_B),
  glance(fit_C)
) |>
  select(.model, AIC, AICc, BIC, sigma2) |>
  arrange(AICc) |>
  kable(digits = 3, align = "c")


# ============================================================
# 7. Head-to-head forecast accuracy
# ============================================================

fc_A <- fit_A |> forecast(new_data = test)
fc_B <- fit_B |> forecast(new_data = test)
fc_C <- fit_C |> forecast(new_data = test)

# Forecast plots
fc_A |>
  autoplot(combined |> filter(yearmonth >= yearmonth("2020 Jan")), level = c(80, 95)) +
  labs(title = "Model A: Forecast vs. Actual Danish Inflation",
       x = "Month", y = "Inflation (YoY, %)") +
  theme(plot.title = element_text(hjust = 0.5))

fc_B |>
  autoplot(combined |> filter(yearmonth >= yearmonth("2020 Jan")), level = c(80, 95)) +
  labs(title = "Model B: Forecast vs. Actual Danish Inflation",
       x = "Month", y = "Inflation (YoY, %)") +
  theme(plot.title = element_text(hjust = 0.5))

fc_C |>
  autoplot(combined |> filter(yearmonth >= yearmonth("2020 Jan")), level = c(80, 95)) +
  labs(title = "Model C (GETS): Forecast vs. Actual Danish Inflation",
       x = "Month", y = "Inflation (YoY, %)") +
  theme(plot.title = element_text(hjust = 0.5))

# Accuracy table: all models side by side.
bind_rows(
  fc_arima_auto   |> accuracy(test),
  fc_arima_breaks |> accuracy(test),
  fc_A            |> accuracy(test),
  fc_B            |> accuracy(test),
  fc_C            |> accuracy(test)
) |>
  select(.model, RMSE, MAE, MAPE, MASE) |>
  arrange(RMSE) |>
  kable(digits = 3, align = "c")


# ============================================================
# 8. Time Series Cross-Validation (TSCV) — ARIMA baseline only
# ============================================================

# TSCV with step dummies is not feasible: in early expanding windows the
# break dummies are constant (all zero), causing rank deficiency. TSCV is
# therefore applied to the ARIMA baseline only.

combined_cv <- combined |>
  stretch_tsibble(.init = 60, .step = 1)

cat("TSCV folds:", max(combined_cv$.id), "\n")

fit_cv <- combined_cv |>
  model(
    arima_cv = ARIMA(dk_inflation, stepwise = TRUE, approximation = TRUE)
  )

fc_cv <- fit_cv |> forecast(h = 1)

cv_accuracy <- fc_cv |>
  accuracy(combined) |>
  select(.model, RMSE, MAE, MAPE, MASE)

cv_accuracy |> kable(digits = 3, align = "c")

bind_rows(
  fc_arima_auto |> accuracy(test) |>
    mutate(.model = "ARIMA baseline (fixed test set, 24-step)"),
  cv_accuracy |>
    mutate(.model = "ARIMA baseline (TSCV, avg. 1-step)")
) |>
  select(.model, RMSE, MAE, MAPE, MASE) |>
  kable(digits = 3, align = "c")

# ============================================================
# GETS IIS/SIS indikator-plot
# ============================================================

# Udtræk bevarede indikatorer fra gets_fit
retained <- names(coef(gets_fit))

# IIS: individuelle udliggere
iis_obs <- retained[grepl("^iis", retained)] |>
  gsub("iis", "", x = _) |>
  as.integer()

# SIS: permanente niveauskift
sis_obs <- retained[grepl("^sis", retained)] |>
  gsub("sis", "", x = _) |>
  as.integer()

# Kortlæg observationsnumre til datoer via gets_df
# (gets_df er træningsdata med komplet case fra lag 6 og frem)
gets_df_dated <- gets_df |>
  mutate(obs = row_number())

iis_dates <- gets_df_dated |>
  filter(obs %in% iis_obs) |>
  pull(yearmonth) |>
  as.Date()

sis_dates <- gets_df_dated |>
  filter(obs %in% sis_obs) |>
  pull(yearmonth) |>
  as.Date()

# Plot
plot_df <- train |>
  as_tibble() |>
  mutate(dato = as.Date(yearmonth))

iis_df <- plot_df |> filter(dato %in% iis_dates)
sis_df <- plot_df |> filter(dato %in% sis_dates)

ggplot(plot_df, aes(x = dato, y = dk_inflation)) +
  geom_line(colour = "black", linewidth = 0.5) +
  geom_vline(xintercept = sis_dates,
             colour = "steelblue", linetype = "dashed", alpha = 0.7) +
  geom_point(data = iis_df, aes(x = dato, y = dk_inflation),
             colour = "firebrick", size = 2.5, shape = 16) +
  annotate("text", x = as.Date("2005-01-01"), y = 10,
           label = "Rød prik = IIS (udligger)", colour = "firebrick",
           hjust = 0, size = 3) +
  annotate("text", x = as.Date("2005-01-01"), y = 9.2,
           label = "Blå linje = SIS (niveauskift)", colour = "steelblue",
           hjust = 0, size = 3) +
  labs(
    title = "GETS: Identificerede udliggere (IIS) og niveauskift (SIS)",
    x     = "Måned",
    y     = "Dansk inflation (ÅoÅ, %)"
  ) +
  theme(plot.title = element_text(hjust = 0.5))

# ============================================================
# GETS: original serie vs. "renset" serie
# ============================================================

# Udtræk koefficienter for IIS og SIS
gets_coefs <- coef(gets_fit)
iis_coefs  <- gets_coefs[grepl("^iis", names(gets_coefs))]
sis_coefs  <- gets_coefs[grepl("^sis", names(gets_coefs))]

# Observation-numre
iis_obs <- as.integer(gsub("iis", "", names(iis_coefs)))
sis_obs  <- as.integer(gsub("sis", "", names(sis_coefs)))

# gets_df har n_tr rækker — byg adjustment-vektor
adjustment <- rep(0, n_tr)

# IIS: tilføj koefficient til den specifikke observation
for (i in seq_along(iis_obs)) {
  obs <- iis_obs[i]
  if (obs <= n_tr) adjustment[obs] <- adjustment[obs] + iis_coefs[i]
}

# SIS: tilføj koefficient fra brudpunkt og frem
for (i in seq_along(sis_obs)) {
  obs <- sis_obs[i]
  if (obs <= n_tr) adjustment[obs:n_tr] <- adjustment[obs:n_tr] + sis_coefs[i]
}

# Byg plot-dataframe
plot_gets <- gets_df |>
  mutate(
    obs        = row_number(),
    dk_infl    = y_tr,
    adjustment = adjustment,
    dk_renset  = dk_infl - adjustment
  ) |>
  select(yearmonth, dk_inflation = dk_infl, dk_renset, adjustment)

plot_gets |>
  pivot_longer(
    cols      = c(dk_inflation, dk_renset),
    names_to  = "serie",
    values_to = "vaerdi"
  ) |>
  mutate(serie = recode(serie,
                        dk_inflation = "Original",
                        dk_renset    = "Renset for IIS/SIS"
  )) |>
  ggplot(aes(x = as.Date(yearmonth), y = vaerdi, colour = serie)) +
  geom_line(linewidth = 0.6) +
  scale_colour_manual(values = c(
    "Original"           = "black",
    "Renset for IIS/SIS" = "steelblue"
  )) +
  labs(
    title  = "Dansk inflation: original vs. renset for IIS/SIS-indikatorer",
    x      = "Måned",
    y      = "Inflation (ÅoÅ, %)",
    colour = NULL
  ) +
  theme(
    plot.title      = element_text(hjust = 0.5),
    legend.position = "bottom"
  )
